% Chapter 3.3: Global Weather Forecasting and Climate Dynamics
% Part III of the book.

\chapter{Global Weather Forecasting and Climate Dynamics}
\label{ch:part3-global-weather-forecasting-and-climate-dyn}
\label{sec:climate}
% -----------------------------------------------------------------------------

Global weather prediction is an initial-value problem on the spherical shell
$\Omega = \mathbb{S}^2 \times [0, z_{\text{top}}]$, and it differs from the aerodynamic
setting of \chapref{ch:part3-external-aerodynamics-and-surface-flows} in the one respect
that governs the whole error budget: nothing about its domain varies. The geometry is a
sphere of fixed radius, so the lattice and stationarity assumptions that
\chapref{ch:part1-the-fourier-neural-operator-spectral-theor} had to surrender for a car
body are available here again, and the transform machinery of Part~II applies without an
encoder or a decoder in front of it. What has become hard is the number of applications.
A ten-day forecast is forty successive uses of a six-hour map; an operational center
integrates an ensemble of some fifty members at each forecast cycle --- the fifty-member,
ten-day benchmark forecast that \cite{pathak2022fourcastnet} takes as its unit of cost
comparison --- and that is the multi-query budget
\chapref{ch:part3-the-multi-query-bottleneck-in-computationa} prices and which
\chapref{ch:part2-production-deployment-and-high-throughput-} notes is already spent
before any surrogate exists. \index{weather forecasting!as an operator problem}
\index{numerical weather prediction (NWP)}

The atmosphere also breaks the error accounting that Parts~I and~II assume. It is chaotic,
so the amplification factor $\Lambda_\theta$ defined in \eqref{eq:ch09-onesided} exceeds one
by an amount the leading Lyapunov exponents fix, the pathwise bound \eqref{eq:ch09-gronwall}
saturates at that rate, and two trajectories with nearly identical starting states separate
while sharing the same statistics; skill is therefore scored field-by-field at a lead
time, not along a trajectory. The labels are not solver output either. They are a
\emph{reanalysis} --- the optimal estimate of the observed history produced by cycling a
data-assimilation system through the observational record \cite{hersbach2020era5} --- so whatever the assimilated
model got wrong is inherited by the surrogate as ground truth. And the quantity a user
cares about is often the part of the field the loss suppresses: the tail of the wind-speed
distribution, the heaviest millimeters of rain, the mesoscale structure that an
area-integral skill score barely notices.

This chapter takes the four questions in order. Section~\ref{sec:ch33-primitive} states the
hydrostatic primitive equations, isolates the three structural features that decide the
design of anything that learns from them, and names the closure and the boundary data the
equations do not supply themselves. Section~\ref{sec:ch33-fourcastnet} builds the actual
surrogate, \textbf{FourCastNet} and its Adaptive Fourier Neural Operator backbone, from the
architecture of Part~II outward: data, state tensor, mixing layer, training procedure.
Section~\ref{sec:ch33-verification} reads the reported skill and cost numbers with their
provenance pinned down, including the distinction between a wall-clock ratio and an
aggregate node-time ratio. Section~\ref{sec:ch33-limits} collects the failures that the same sources
report --- smoothed small scales, underestimated extremes, no extrapolation to unseen
intensity --- and states plainly why none of them licenses a climate claim. The reader
should bring the mode budget \eqref{eq:ch05-modeset}, the rollout recursion
\eqref{eq:ch26-rollout}, and the pushforward objective \eqref{eq:ch26-pushforward}; the
chapter adds a domain, a data policy, and a verification metric, not new theory.

% -----------------------------------------------------------------------------
\section{Governing Hydrodynamic Primitive Equations}
\label{sec:ch33-primitive}
% -----------------------------------------------------------------------------

Atmospheric circulation at the scales that matter to a forecast is modeled by the
hydrostatic primitive equations on a rotating sphere:
\begin{align}
  \label{eq:ch33-momentum}
  \frac{D \mathbf{u}_h}{Dt} + f \hat{\mathbf{k}} \times \mathbf{u}_h &= -\frac{1}{\rho} \nabla_h p + \mathbf{F}_h, \\
  \label{eq:ch33-hydrostatic}
  \frac{\partial p}{\partial z} &= -\rho g, \\
  \label{eq:ch33-continuity}
  \nabla_h \cdot \mathbf{u}_h + \frac{\partial w}{\partial z} &= 0, \\
  \label{eq:ch33-pottemp}
  \frac{D \theta}{Dt} &= \frac{Q}{c_p} \left(\frac{p_0}{p}\right)^{R/c_p},
\end{align}
where $\mathbf{u}_h = (u, v)$ is horizontal velocity, $f = 2\Omega_{\text{rot}} \sin\phi$ is
the Coriolis parameter at latitude $\phi$ for planetary rotation rate $\Omega_{\text{rot}}$,
and $\theta$ is potential temperature. The remaining symbols are the horizontal gradient
$\nabla_h$ and divergence on $\mathbb{S}^2$, the local vertical unit vector
$\hat{\mathbf{k}}$, density $\rho$, pressure $p$, vertical velocity $w$, gravitational
acceleration $g$, the material derivative
$D/Dt = \partial_t + \mathbf{u}_h \cdot \nabla_h + w\,\partial_z$, specific heat at
constant pressure $c_p$, the gas constant for dry air $R$, a reference pressure $p_0$, and
the two forcing terms $\mathbf{F}_h$ and $Q$. The derivation and the scaling that justifies
each term are standard geophysical fluid dynamics \cite{vallis2017atmospheric}, and
\eqref{eq:ch33-momentum}--\eqref{eq:ch33-pottemp} are used here as the definition of what the
surrogate is being asked to imitate. \index{primitive equations!hydrostatic}
\index{Coriolis parameter}

Three structural features of \eqref{eq:ch33-momentum}--\eqref{eq:ch33-pottemp} decide the
design of anything that learns from them. The first is balance. Away from acoustic and
gravity-wave noise, \eqref{eq:ch33-momentum} is close to geostrophic balance,
$f\hat{\mathbf{k}}\times\mathbf{u}_h \approx -\rho^{-1}\nabla_h p$, and
\eqref{eq:ch33-hydrostatic} is hydrostatic balance, so the state is not an arbitrary
element of a high-dimensional space: the fast modes carry little of the information a
verification metric rewards. This is why the geopotential height at $500\ \mathrm{hPa}$,
a smooth mid-tropospheric proxy for the pressure field, is the conventional index of
forecast quality in Section~\ref{sec:ch33-verification}, and why a surrogate that nails it
can still be wrong about the surface winds a wind farm needs. The second is that
\eqref{eq:ch33-continuity} is a constraint rather than an evolution equation: it makes the
vertical velocity diagnostic, obtained by integrating the horizontal divergence, and it
holds the mass of each atmospheric column in balance with the flux through its walls.
Nothing in \eqref{eq:ch33-statemap} enforces it, so a learned state may satisfy
\eqref{eq:ch33-momentum} and \eqref{eq:ch33-pottemp} to within its loss while quietly
creating or destroying column mass --- the same defect
\chapref{ch:part3-boundary-condition-enforcement-neural-oper} prices for the aerodynamic
surrogate, but there invisible in a single field and here accumulating over forty steps.
The third is that \eqref{eq:ch33-momentum} and \eqref{eq:ch33-pottemp} are not closed.
$\mathbf{F}_h$ lumps subgrid momentum mixing and boundary-layer turbulence; $Q$ lumps
radiative transfer and the latent heat release of convection, both of which are driven by
processes orders of magnitude smaller than the $28\ \mathrm{km}$ cell of
Section~\ref{sec:ch33-fourcastnet} and are supplied by parameterization schemes rather than
by the equations themselves.
\index{geostrophic balance}\index{physical parameterization!convection and radiation}

The system is also incomplete in the ordinary boundary-value sense, and the omissions
matter for what can be learned. \eqref{eq:ch33-momentum}--\eqref{eq:ch33-pottemp} require an
initial condition on the shell plus side conditions: no flow through the ground, vanishing
vertical velocity at the model top, periodicity in longitude, regularity at the poles. The
vertical coordinate implicit in $z$ is not geometric height but a hybrid coordinate,
pressure-like aloft and conforming to orography near the surface, so the same symbol denotes
a different physical spacing at every level and the grid is baked into the archive the
surrogate reads. The state that determines the future is larger still than the five fields
displayed: surface temperature over land, sea ice, and ocean, soil moisture, and the
radiative environment all enter $\mathbf{F}_h$ and $Q$, which is why a working surrogate is
handed twenty coupled channels rather than the primitive variables alone.

The consequence for operator learning is worth stating exactly. A solution operator is
defined by a differential equation together with its closure; here the closure is a
collection of tuned, non-exact models whose errors are a leading contributor to precipitation
forecast error. So the map an atmospheric surrogate learns is not a partial
differential equation (PDE) $\to$ solution map in the sense of
\chapref{ch:part1-partial-differential-equations-as-solution}. It is
$\text{PDE} + \text{parameterization package} + \text{assimilation system} \to \text{analysis sequence}$,
and every later claim of ``accuracy'' in this chapter is accuracy against that composite.
A classical solver at least knows which of its terms produced an error. A single learned
map does not, and cannot be audited term by term --- an instance of the opacity that
\chapref{ch:part3-boundary-condition-enforcement-neural-oper} treats from the
conservation side.

% -----------------------------------------------------------------------------
\section{FourCastNet and Adaptive Fourier Neural Operators}
\label{sec:ch33-fourcastnet}
% -----------------------------------------------------------------------------

\textbf{FourCastNet}~\cite{pathak2022fourcastnet} parameterizes the atmospheric solution
operator using \textbf{Adaptive Fourier Neural Operators} (AFNO), the spectral token mixer
of Guibas et al.~\cite{guibas2021afno}, whose construction the FourCastNet paper sets out
alongside the model that uses it. Operating on
the spherical equispaced latitude--longitude grid from the \textbf{European Centre for
Medium-Range Weather Forecasts} (ECMWF) \textbf{ERA5} reanalysis dataset
\cite{pathak2022fourcastnet,hersbach2020era5} --- $720 \times 1440$ pixels, corresponding to
$0.25^\circ \approx 28 \text{ km}$ spatial resolution --- the model predicts 20 prognostic
atmospheric variables:
\begin{equation}
  \label{eq:ch33-statemap}
  \mathcal{G}_{\text{AFNO}} : \mathbf{X}(t) \in \mathbb{R}^{20 \times 720 \times 1440} \longmapsto \mathbf{X}(t + 6\text{ h}).
\end{equation}
Equation \eqref{eq:ch33-statemap} is the time-$\Delta t$ map of
\chapref{ch:part1-open-theoretical-problems-and-mathematical} with $\Delta t = 6\ \mathrm{h}$
and the state $\mathbf{X}$ a twenty-channel image on the sphere; a $K$-step forecast is
$\mathcal{G}_{\text{AFNO}}^{K}$ applied as in \eqref{eq:ch26-rollout}. Two bookkeeping
details matter for reproduction. The archived regular grid carries $721$ latitude nodes
because both poles are included, so the tensor the code actually reads is
$721 \times 1440 \times 20$; and ERA5's native horizontal grid is a \emph{Gaussian} grid,
re-interpolated to the regular latitude--longitude arrangement by the Copernicus Climate
Data Store service before training \cite{pathak2022fourcastnet}. The interpolation is
benign arithmetically and not benign physically: it is a low-pass operation applied to the
labels, and any claim about what the model resolves must be read against it. \index{ERA5 reanalysis}\index{FourCastNet}
\index{latitude--longitude grid!pole singularity}

The geometry of that grid is the reason the operator's global mixing is cheap and the
reason its fine scales are not isotropic. At $0.25^\circ$ in both directions, a cell is
$28\ \mathrm{km}$ wide in the meridional sense everywhere, while its zonal width shrinks as
$\cos\phi$: the row one step below the pole is roughly $120\ \mathrm{m}$ wide in longitude
while its meridional extent stays $28\ \mathrm{km}$, an aspect ratio near $2\times10^{2}$, and
the polar row itself --- the $721$st latitude node --- is $1440$ stored copies of a single
point.
A spectral layer on such an array is not a spectral layer on the
sphere, since the spherical Laplacian's eigenfunctions are spherical harmonics, whose
degree-$n$ energy sits on cells of increasingly anisotropic shape. Operational models
handle this with reduced grids or with spectral-transform methods on Gaussian grids; the
surrogate instead treats the field as a Euclidean image and inherits the anisotropy. This
is a genuine departure from the periodic-domain setting of
\chapref{ch:part1-the-fourier-neural-operator-spectral-theor}, where the lattice is uniform
by assumption, and it is one reason the small-scale results of
Section~\ref{sec:ch33-limits} come out as they do. \index{spherical harmonics!versus latitude--longitude grids}

The channels in \eqref{eq:ch33-statemap} are chosen to represent the instantaneous state
rather than to be complete. Surface wind components $U_{10}, V_{10}$ at $10\ \mathrm{m}$,
2-meter temperature $T_{2m}$, surface pressure and mean sea-level pressure; zonal and
meridional wind and geopotential height at $1000\ \mathrm{hPa}$; temperature, wind,
geopotential and relative humidity at $850$ and $500\ \mathrm{hPa}$; geopotential height at
$50\ \mathrm{hPa}$; and integrated total column water vapor (TCWV) --- twenty fields in
all. Total precipitation is deliberately \emph{not} among them: it is a six-hour
accumulated diagnostic, and a second AFNO model maps the backbone's output onto it, with a
final convolution and rectified linear unit appended so that negative rainfall cannot be
produced \cite{pathak2022fourcastnet}.
\index{precipitation!as a diagnostic head}

Inside, the mixer is the AFNO layer rather than the \textbf{Fourier Neural Operator} (FNO)
layer of \eqref{eq:fno_layer}. The $720 \times 1440$ frame is projected onto an $h \times w$
grid of patches of size $p \times p$, with $p = 8$ in the published configuration, each patch
embedded as a $d$-dimensional token; the
tokens are transformed along the patch indices by a two-dimensional discrete Fourier
transform, a two-layer perceptron with block-diagonal weights shared across all patches acts
on each frequency's token embedding, and the result is sparsified and returned to the patch
domain with a residual connection \cite{pathak2022fourcastnet,guibas2021afno}. The sparsification is the
adaptive part and is what distinguishes it from a fixed low-pass truncation:
\begin{equation}
  \label{eq:ch33-shrink}
  \tilde{z}_{m,n} = S_\lambda\!\left(\mathrm{MLP}\left(z_{m,n}\right)\right),
  \qquad
  S_\lambda(x) = \mathrm{sign}(x)\,\max\!\left(|x| - \lambda,\, 0\right),
\end{equation}
where $z_{m,n}$ is the token embedding at frequency index $(m,n)$, $\mathrm{MLP}$ is the
two-layer perceptron with block-diagonal weights just described, and $\lambda$ is the
sparsity-controlling threshold, listed as $10^{-2}$ in the published hyperparameters. Equation
\eqref{eq:ch33-shrink} retains the modes whose coefficients survive
shrinkage instead of retaining a box such as \eqref{eq:ch05-modeset} by construction, so
the frequency support is input- and channel-dependent rather than fixed; the number of
parameters does not scale with the number of patches, which is the resolution-invariance of
\chapref{ch:part1-discretization-invariance-and-continuum-li} as an architectural property
of the mixer. \index{adaptive Fourier neural operator (AFNO)!token mixing}
\index{soft-thresholding!in the spectral domain}

The engineering payoff of that choice is stated in the source paper in memory terms, and it
is one of the two reasons that paper gives for preferring a transformer-style backbone over a
convolutional one: at \eqref{eq:ch33-statemap} resolution the AFNO model's footprint is about
$10\ \mathrm{GB}$ at batch size one, whereas carrying a nineteen-layer convolutional network
previously trained at $32 \times 64$ over to the same frame would need roughly
$83\ \mathrm{GB}$, and its receptive field would have to grow with the grid. The other reason
is empirical rather than arithmetic: the authors report that their convolutional runs
``showed poor performance on capturing small scales over many time steps in auto-regressive
inference'' \cite{pathak2022fourcastnet}, which is why Section~\ref{sec:ch33-limits} has a
subject at all. Compare the spectral-layer costing of \eqref{eq:ch23-cost}: the
$\mathcal{O}(N\log N)$ mixing is the same object, but the alternative to paying it is not a
cheaper network, it is a deeper one.

Training uses hourly ERA5 snapshots subsampled to $00, 06, 12, 18\ \mathrm{UTC}$, with
$54{,}020$ training frames drawn from 1979--2015, $2{,}920$ validation frames from
2016--2017, and 2018 onward held out untouched as the out-of-sample test
set \cite{pathak2022fourcastnet}. The procedure is two-stage and should be
recognized from Part~II: a one-step pretraining on pairs
$\big(\mathbf{X}(k),\mathbf{X}(k+1)\big)$, then fine-tuning in which the model is rolled out
two steps on its own output and the loss is the sum of the two step errors. That is
pushforward training \eqref{eq:ch26-pushforward} in all but name, at $K=2$. The source states
the procedure without justifying it;
\chapref{ch:part2-training-regimes-optimization-and-failure-} supplies the justification,
which is exposure bias under the rollout recursion \eqref{eq:ch26-rollout}. End-to-end
training takes about sixteen hours on sixty-four NVIDIA A100 graphics processing units
(GPUs).
\index{pushforward regularization!two-step fine-tuning}

% -----------------------------------------------------------------------------
\section{Comparison with Classical Operational Weather Models (ECMWF IFS)}
\label{sec:ch33-verification}
% -----------------------------------------------------------------------------

The baseline is \textbf{numerical weather prediction} (NWP) with the ECMWF
\textbf{Integrated Forecasting System} (IFS), a spectral-transform dynamical core solving
\eqref{eq:ch33-momentum}--\eqref{eq:ch33-pottemp} with parameterized physics on a CPU
cluster. The comparison is made on the metric operational services already use. For a
forecast $\mathbf{X}$, a verification analysis $\mathbf{Y}$, and climatology
$\mathbf{C}$, the \textbf{Anomaly Correlation Coefficient} (ACC) is
\begin{equation}
  \label{eq:ch33-acc}
  \text{ACC} = \frac{\sum_{i,j} (\mathbf{X}_{ij} - \mathbf{C}_{ij}) (\mathbf{Y}_{ij} - \mathbf{C}_{ij})}{\sqrt{\sum_{i,j} (\mathbf{X}_{ij} - \mathbf{C}_{ij})^2 \sum_{i,j} (\mathbf{Y}_{ij} - \mathbf{C}_{ij})^2}},
\end{equation}
where the anomaly is the departure from a climatological mean, so that
\eqref{eq:ch33-acc} is blind to any additive bias in $\mathbf{X}$: a forecast displaced
uniformly from the analysis still scores one. The \textbf{root-mean-square error} (RMSE)
\index{root-mean-square error (RMSE)}
reported alongside it is the quantity such a displacement would move, which is why the two
are always quoted together. The conventional instance of \eqref{eq:ch33-acc} is the
geopotential height at $500\ \mathrm{hPa}$, one of the twenty channels of
\eqref{eq:ch33-statemap}: smooth, mid-tropospheric, and the standard proxy for the pressure
field aloft.
\index{anomaly correlation coefficient (ACC)}

Equation \eqref{eq:ch33-acc} as written is a sum over array indices, and on a sphere that
is a defect rather than a simplification: the same field scored on a finer longitude
spacing would receive a different number, because the sum weights a square kilometer near
the pole far more heavily than a square kilometer near the equator. This is the
grid-dependence that \chapref{ch:part1-discretization-invariance-and-continuum-li}
identifies in the loss, and the remedy is the same --- insert the quadrature weights of the
area element $\mathrm{d}A = R^2 \cos\phi\, \mathrm{d}\phi\, \mathrm{d}\lambda$:
\begin{equation}
  \label{eq:ch33-accw}
  \text{ACC}_w = \frac{\sum_{i,j} w_j (\mathbf{X}_{ij} - \mathbf{C}_{ij}) (\mathbf{Y}_{ij} - \mathbf{C}_{ij})}{\left(\sum_{i,j} w_j (\mathbf{X}_{ij} - \mathbf{C}_{ij})^2 \; \sum_{i,j} w_j (\mathbf{Y}_{ij} - \mathbf{C}_{ij})^2\right)^{1/2}},
  \qquad w_j = \cos\phi_j .
\end{equation}
The published FourCastNet verification uses \eqref{eq:ch33-accw}, together with the
correspondingly weighted RMSE, averaged over many
initial conditions spread through a full year so that seasonal variation in skill does not
enter as a bias \cite{pathak2022fourcastnet}. \index{area weighting!on the sphere}

What the numbers say, with the hedging the source itself supplies. ACC is a correlation
that decays with lead time, so a forecast horizon is whatever threshold one agrees to
cross; the FourCastNet authors report ACC above $0.6$ for TCWV beyond day eight, and state
that their predictions are comparable to the IFS on ACC and RMSE at lead times of up to
three days, after which ``predictions of all modeled variables lag
close behind IFS'' out to a week \cite{pathak2022fourcastnet}. The stronger version of the
claim --- that the surrogate matches the operational solver at short lead times for
large-scale variables while outperforming it for small-scale variables, precipitation among
them --- is also the authors' own, and it is a claim
about $\mathbf{X}$ against ERA5 at a fixed lead time, not about the statistics of a long
rollout. \index{forecast horizon!defined by a skill threshold}

Cost accounting requires the same care, because the ratio depends entirely on what is being
divided. The paper's own comparison is node-to-node for a $100$-member, $24$-hour ensemble:
the IFS at $18\ \mathrm{km}$ is estimated to need roughly $984{,}000$ node-seconds and about
$271\ \mathrm{MJ}$, while FourCastNet completes the same task in about $7$ node-seconds on
a single four-GPU node at roughly $8\ \mathrm{kJ}$, the IFS needing on the order of
$3{,}060$ nodes for the job. The resulting speedups are $145{,}000\times$ at the
$30\ \mathrm{km}$ the trained model runs at --- the paper's rounding of the
$0.25^\circ$ grid --- and $45{,}000\times$ for an $18\ \mathrm{km}$ realization of it, with
energy ratios of $24{,}000\times$ and $12{,}000\times$ respectively; the second of those two
speedups is an extrapolation rather than a
measurement, since no $18\ \mathrm{km}$ version of the model was ever trained and its
timing is obtained by interpolating data and parameters to that grid
\cite{pathak2022fourcastnet}. A second, far smaller ratio is the one usually repeated: a
week-long global forecast in under two seconds on one A100 --- $1.8\ \mathrm{s}$, the
figure this book carries; this book also quotes roughly two hours of
wall-clock for a single IFS member on a supercomputer cluster. Their ratio,
$7{,}200/1.8$, is the familiar $4{,}000\times$. That division is arithmetic on two different
machines, and the primary source does not report the elapsed time of one member on one
partition at all: its only wall-clock figure is for a whole ensemble, a $15$-day,
$51$-member IFS run on $1{,}530$ CPU nodes in $82$ minutes, from which the node-second
totals above are derived. This chapter therefore reads the $4{,}000\times$ as a statement
about latency on a device rather than as a costing. The two accountings answer different
questions, and the
difference between $4{,}000\times$ and $45{,}000\times$ is mostly the ensemble size divided
out. Table~\ref{tab:ch33-cost} sets them side by side. \index{speedup!wall-clock versus node-hour}
\index{European Centre for Medium-Range Weather Forecasts (ECMWF)}
\index{Integrated Forecasting System (IFS)}

\begin{table}[htbp]
\centering
\small
\caption{Two accountings of the same surrogate. Wall-clock ratios compare one forecast
member on one device; aggregate ratios compare the machine time of a whole ensemble.
Values as reported or estimated by \cite{pathak2022fourcastnet}; the per-member wall-clock
figures in the third row are the ones this book carries.}
\label{tab:ch33-cost}
\begin{tabular}{@{}>{\raggedright\arraybackslash}p{0.30\textwidth}>{\raggedright\arraybackslash}p{0.30\textwidth}>{\raggedright\arraybackslash}p{0.32\textwidth}@{}}
\toprule
\textbf{Quantity} & \textbf{ECMWF IFS} & \textbf{FourCastNet (AFNO)} \\
\midrule
State represented & $>150$ variables, $>50$ vertical levels & 20 variables at five levels, surface included \\
Time stepping & spectral-transform dynamical core and parameterized physics & learned six-hour map \eqref{eq:ch33-statemap} \\
Week-long global forecast & $\approx 2\ \mathrm{h}$ per member; $82$ min for a $15$-day $51$-member run on $1{,}530$ nodes & under $2\ \mathrm{s}$ ($1.8\ \mathrm{s}$ as quoted) on one A100 \\
$100$-member, $24$-h ensemble & $\approx 984{,}000$ node-s, $\approx 271\ \mathrm{MJ}$, $\approx 3{,}060$ nodes & $\approx 7$ node-s, $\approx 8\ \mathrm{kJ}$, one $4\times$A100 node \\
Aggregate speedup & --- & $45{,}000\times$ ($18\ \mathrm{km}$, extrapolated) to $145{,}000\times$ ($30\ \mathrm{km}$) \\
Energy ratio & --- & $12{,}000\times$ to $24{,}000\times$ lower \\
Development cost & decades of model development & $\approx 16\ \mathrm{h}$ on $64$ A100 GPUs; training energy $\approx$ one 10-day 50-member IFS forecast \\
\bottomrule
\end{tabular}
\end{table}

Table~\ref{tab:ch33-cost} makes the real argument of this chapter visible in one place: the
surrogate's advantage is not that it forecasts, which the IFS also does, but that it makes
ensemble size a free variable. The IFS reference run behind the accounting above is a
$51$-member ensemble, and every extra member costs another slice of a $1{,}530$-node
allocation; a thousand-member ensemble of a $7$ node-second model, by contrast, is an
accounting exercise rather than a project, and initial-condition perturbation by Gaussian
noise on the standardized state is enough to produce one \cite{pathak2022fourcastnet}.
Part~II's
throughput machinery of \chapref{ch:part2-production-deployment-and-high-throughput-} is
exactly what turns that arithmetic into a batch dimension. Panel~(b) of
Figure~\ref{fig:vol3_aero}, printed in \chapref{ch:part3-external-aerodynamics-and-surface-flows},
shows the object these numbers are attached to: the sphere. That plate's other panel,
(a), carries the vehicle surface treated in the present chapter; the two are printed
together there because they are the same transform applied to a domain that is fixed
and to one that is not.
\index{ensemble forecasting!cost of}


% -----------------------------------------------------------------------------
\section{Limits of the Operator Reading: Scale, Extremes, and Climate}
\label{sec:ch33-limits}
% -----------------------------------------------------------------------------

Panel~(b) of Figure~\ref{fig:vol3_aero} is also the first clue to what is missing. Its vorticity
contours are smooth and its jets are broad; nothing in that panel shows a squall line, a
frontal band, or the sub-structure of a precipitation field. The quantitative literature
agrees with the eye test. Bonavita's 2024 audit~\cite{bonavita2024limitations}, which examines FourCastNet alongside two
other leading learned forecasters, concludes that such models ``are not
able to properly reproduce sub-synoptic and mesoscale weather phenomena and lack the
fidelity and physical consistency of physics-based models,'' with consequences for the
interpretation of their forecasts and of their perceived skill. The mechanism is not
mysterious. A least-squares objective against a single future state estimates the conditional
mean of
that state given the input; the conditional mean of a field with sharp fronts is a smoothed
field, and repeated application of \eqref{eq:ch33-statemap} compounds the smoothing. This
is the same arithmetic as the coefficient-decay bound \eqref{eq:ch05-coefficient-decay},
with the learned shrinkage of \eqref{eq:ch33-shrink} in place of a hard mode box.
\index{mesoscale!under-resolved by learned forecasters}\index{conditional mean!smoothing toward}

The model's own authors measure the consequence for tails rather than asserting it. They
define a \textbf{relative quantile error} (RQE) over logarithmically spaced upper quantiles
$Q = \{1-10^{-1}, \dots, 1-10^{-4}\}$,
\begin{equation}
  \label{eq:ch33-rqe}
  \mathrm{RQE}(\ell) = \sum_{q \in Q} \frac{\mathbf{X}^{q}_{\text{pred}}(\ell) - \mathbf{X}^{q}_{\text{true}}(\ell)}{\mathbf{X}^{q}_{\text{true}}(\ell)},
\end{equation}
where $\mathbf{X}^{q}(\ell)$ is the $q$-th quantile of the field at lead-time index $\ell$,
so that a negative value records systematic under-prediction of extremes. Equation
\eqref{eq:ch33-rqe} trends negative for both the surrogate and the IFS on surface wind
speed --- a few percentage points, with negligible difference between them --- and much
more strongly for precipitation, where the reported under-prediction of the extreme
percentiles is roughly $35\%$ against roughly $15\%$ for the IFS
\cite{pathak2022fourcastnet}. The same paper reports that a $100$-member ensemble tracks
Hurricane Michael's position well over its four-day trajectory but fails to reproduce the
full extent of the sharp drop in minimum sea-level pressure between $36$ and $48\ \mathrm{h}$,
and its authors hypothesize
that the failure comes from processes the model does not represent, convective and radiative
among them. That is the
closure of Section~\ref{sec:ch33-primitive} reappearing as a forecast error. \index{relative quantile error (RQE)}
\index{extremes!under-prediction of}

Out-of-distribution intensity is a harder limit than under-prediction, because it is not a
bias one can correct. Sun et al., in a 2025 study~\cite{sun2025grayswan}, retrained FourCastNet four times
on the $1979$--$2015$ ERA5 record: once on all of it, and three times with Category~3--5
tropical cyclones removed either globally or over a single ocean basin. Tested on the
Category~5 cyclones of $2018$--$2023$, the versions give essentially unchanged global
verification skill, but the version trained without Category~3--5 storms anywhere cannot
forecast them: the model does not extrapolate from weaker storms present in the
data to stronger ones absent from it. For a hazard-planning user, this is the property that
matters, and no ACC reported at day five speaks to it. \index{out-of-distribution!gray-swan cyclones}

Finally, the distinction this chapter's title invites and its evidence does not support.
\eqref{eq:ch33-statemap} is fit to short-horizon transitions between analyzed states inside
a record spanning four decades of one forced realization of the climate system. Climate
questions are questions about the invariant measure of the dynamics and about its response
to a forcing the training record does not contain. Part~I's rollout analysis is explicit
that nothing licenses the extrapolation: the pathwise bound \eqref{eq:ch09-gronwall}
becomes vacuous once the Lyapunov growth is resolved, the useful requirement is an
absorbing or statistical one such as \eqref{eq:ch09-weak}, and a spatial correlation at one
lead time says nothing about whether $\mathcal{G}_{\theta}^{K}$ has the right attractor as
$K \to \infty$. A model that under-represents mesoscale
variance, as \eqref{eq:ch33-rqe} shows, will misrepresent the statistics of any quantity
that variance sets --- eddy heat transport, storm-track frequency, the intensity
distribution --- even while its day-three anomaly correlations are excellent. Data
availability compounds this: ERA5 is a product of an assimilation system that changed
across the record \cite{hersbach2020era5}, so a learned map trained on the whole archive fits a non-stationary
estimator as well as a non-stationary atmosphere. \index{weather forecasting!weather versus climate}
\index{climate emulation!not provided by weather surrogates}
\index{reanalysis!as training labels}

\takeaway{What the Weather Numbers License}{A learned six-hour state map on the sphere
reaches IFS-comparable anomaly correlation for large-scale variables out to about three
days, lags slightly behind thereafter, and makes a $100$-member daily ensemble roughly four
to five orders of magnitude cheaper in aggregate machine time. It does not resolve
sub-synoptic structure, it under-predicts precipitation extremes by about $35\%$, it cannot
extrapolate to cyclone intensities absent from its training record, and none of its skill
scores constitutes evidence about climate.}

% -----------------------------------------------------------------------------
% Bridge to Chapter 3.4
% -----------------------------------------------------------------------------

The atmospheric case is the cleanest in this book because its geometry never changes, and
that is precisely what makes its failures diagnostic: smoothing, tail under-prediction, and
distribution shift are properties of the learned map and the loss, not of an awkward
domain. Move the same three failures onto a domain where the geometry does change and where
the front is not merely under-resolved but genuinely discontinuous, and the difficulty
sharpens. \chapref{ch:part3-subsurface-multiphase-flow-and-carbon-stor} does exactly that.
There the operator maps a permeability field and an injection schedule to the saturation of
a supercritical phase --- lighter than the brine it displaces, which is why it rises ---
migrating through layered rock beneath a low-permeability seal,
over decades rather than days; the governing two-phase Darcy system has moving sharp
saturation fronts whose spectral coefficients decay like the slow tail of
\eqref{eq:ch05-coefficient-decay}, so a fixed mode budget is a worse approximation than the
one the atmosphere tolerates. The remedy taken up there is architectural rather than
statistical --- multi-resolution branches that restore the local detail a global transform
suppresses --- and the verification target changes with it, from a correlation against an
analysis to plume extent, front position, and pressure buildup at a well. The data regime
changes too: where the atmosphere supplied four decades of hourly global analyses, the
reservoir supplies a handful of core measurements and a stochastic geological prior, so the
multi-query budget of \chapref{ch:part3-the-multi-query-bottleneck-in-computationa} returns
as uncertainty quantification over permeability fields rather than as an ensemble of initial
conditions.

% -----------------------------------------------------------------------------
